\begin{theorem}[Stieltjes 刚性：行列式势与 Schatten 谱完备表征纤维谱]\label{thm:pom-fiber-spectrum-stieltjes-rigidity-determinant-schatten}
\leanverified{paper\_pom\_fiber\_spectrum\_stieltjes\_rigidity\_determinant\_schatten}
固定分辨率 \(m\)，记纤维多重度 \(d_m(x):=\card{\Fold_m^{-1}(x)}\)（\(x\in X_m\)）。
令对角算子 \(\mathcal F_m\) 作用在 \(\ell^2(X_m)\) 的标准基上为
\(
(\mathcal F_m e_x):=d_m(x)\,e_x
\)，并定义行列式势
\[
Z_m(t):=\det\!\bigl(I+t\mathcal F_m\bigr)=\prod_{x\in X_m}\bigl(1+t\,d_m(x)\bigr).
\]
令
\[
\mathcal S_m(t):=\frac{\dd}{\dd t}\log Z_m(t)=\sum_{x\in X_m}\frac{d_m(x)}{1+t\,d_m(x)}.
\]
则成立：
\begin{enumerate}
  \item \(\mathcal S_m(t)\) 在 \(t>0\) 上为 Stieltjes 型函数；其作为有理函数的全部极点恰为
  \(
  t=-1/d
  \)
  （\(d\) 遍历多重集 \(\{d_m(x)\}_{x\in X_m}\)），并且在 \(t=-1/d\) 处的留数等于该 \(d\) 在谱中的出现重数。
  \item 因而多重集 \(\{d_m(x)\}_{x\in X_m}\) 可由 \(\mathcal S_m\) 的极点--留数结构唯一重建；等价地，\(Z_m(t)\) 的零点多重集为 \(\{-1/d_m(x)\}_{x\in X_m}\)。
  \item 若另外给定条件期望 \(E_m\) 的偶阶 Schatten 数据 \(\{\|E_m\|_{2k}\}_{k\ge 1}\)，并且其奇异值为 \(\{d_m(x)^{-1/2}\}_{x\in X_m}\)，则
  \[
  \|E_m\|_{2k}^{2k}=\sum_{x\in X_m}d_m(x)^{-k}\qquad(k\ge 1),
  \]
  从而得到所有负幂和 \(\sum_x d_m(x)^{-k}\)。
  由 Newton 恒等式可恢复多项式 \(\prod_x(1+t\,d_m(x))\)，并因此恢复 \(Z_m(t)\) 与纤维谱的全部多重信息。
\end{enumerate}
\end{theorem}

\begin{proof}
由对角化定义直接得到
\(
Z_m(t)=\prod_x(1+t\,d_m(x))
\)，从而
\(
\mathcal S_m(t)=\sum_x \frac{d_m(x)}{1+t\,d_m(x)}
\)。
每一项在 \(t=-1/d_m(x)\) 处为简单极点，且留数为 \(1\)；对同一 \(d\) 的项相加给出留数等于出现重数。
第二条为有理函数的部分分式分解与零点集合的唯一性。
第三条中，若 \(E_m\) 的奇异值为 \(d_m(x)^{-1/2}\)，则 Schatten--\(2k\) 范数的定义给出
\(
\|E_m\|_{2k}^{2k}=\sum_x (d_m(x)^{-1/2})^{2k}=\sum_x d_m(x)^{-k}
\)。
由所有负幂和可经 Newton 恒等式恢复特征多项式的系数，进而恢复其根多重集 \(\{-1/d_m(x)\}\) 与 \(Z_m(t)\)。
证毕。
\end{proof}

\leanverified{paper\_pom\_conditional\_expectation\_schatten\_negative\_pressure}
\begin{corollary}[条件期望的 Schatten 剖面与负压强]\label{cor:pom-conditional-expectation-schatten-negative-pressure}
在定理 \ref{thm:pom-fiber-spectrum-stieltjes-rigidity-determinant-schatten} 的设定下，若条件期望 \(E_m\) 的奇异值多重集为 \(\{d_m(x)^{-1/2}:x\in X_m\}\)，则对任意 \(p\in[1,\infty)\) 有精确恒等式
\[
\boxed{\;
\|E_m\|_{\mathcal S_p}^{\,p}=\sum_{x\in X_m}d_m(x)^{-p/2}=S_{-p/2}(m)\;}
\]
其中 \(S_s(m):=\sum_x d_m(x)^s\)。
从而
\[
\limsup_{m\to\infty}\frac{1}{m}\log\|E_m\|_{\mathcal S_p}
=\frac{1}{p}\,\limsup_{m\to\infty}\frac{1}{m}\log S_{-p/2}(m).
\]
并且由反射不等式（命题 \ref{prop:pom-pressure-reflection-inequality}）推出下界
\[
\limsup_{m\to\infty}\frac{1}{m}\log\|E_m\|_{\mathcal S_p}
\ \ge\ \frac{2\log\varphi}{p}-\frac{1}{p}\,\limsup_{m\to\infty}\frac{1}{m}\log S_{p/2}(m),
\]
其中 \(\varphi=(1+\sqrt5)/2\)。
\end{corollary}
\begin{proof}
第一式由 Schatten--\(p\) 范数定义与奇异值假设直接给出：
\(
\|E_m\|_{\mathcal S_p}^p=\sum_x (d_m(x)^{-1/2})^p=\sum_x d_m(x)^{-p/2}
\)。
取对数并除以 \(m\) 得到第二式。最后一条由命题 \ref{prop:pom-pressure-reflection-inequality} 的
\(
S_{p/2}(m)S_{-p/2}(m)\ge |X_m|^2
\)
取对数并用 \(\lim_{m\to\infty}\frac{1}{m}\log|X_m|=\log\varphi\) 得到。证毕。
\end{proof}

\endinput
